\chapter{Prefix Sums, Two Pointers and Sliding Windows}\label{ch:prefix}

A surprising number of problems ask about \emph{subarrays}: their sums, their distinct values, their medians, the longest
or shortest one with some property. An array of $n$ elements has about $n^2/2$ subarrays, so looking at each is quadratic.
This chapter collects the techniques that bring such questions down to linear or near-linear time: prefix sums, which answer
any range sum in constant time; difference arrays, their inverse; the hash-map trick for counting subarrays with a given
sum; two pointers, which slide a window across the array while touching every element only twice; and the stacks and
multisets that maintain more complex information about a sliding window. Together they solve a large fraction of all Silver
problems and appear as building blocks inside countless harder ones.

\begin{chapterintro}{\bronze\ \silver\ \gold}{Chapters~\ref{ch:complexity}, \ref{ch:ds-basics} and~\ref{ch:sorting}.}
  \item answer range-sum queries in one and two dimensions in constant time;
  \item apply many range updates offline with difference arrays, including rectangles and arithmetic progressions;
  \item count subarrays with a given sum, remainder or balance using prefix sums and a hash map;
  \item recognise when two pointers apply, and prove that they run in linear time;
  \item maintain the median of a sliding window, and find largest rectangles with monotonic stacks.
\end{chapterintro}

% ======================================================================================================================
\section[Prefix sums]{Prefix sums \hfill \bronze}\label{sec:prefix}

Define $P_0 = 0$ and $P_{i+1} = P_i + a_i$, so that $P_i$ is the sum of the first $i$ elements $a_0, \ldots, a_{i-1}$. Then
the sum of any range is the difference of two prefix sums:
\[
  a_l + a_{l+1} + \cdots + a_{r-1} = P_r - P_l ,
\]
because $P_r$ contains exactly the elements $a_0, \ldots, a_{r-1}$ and subtracting $P_l$ removes $a_0, \ldots, a_{l-1}$
(Figure~\ref{fig:prefix}). The array $P$ takes $\Oh{n}$ to build and then answers every range sum in $\Oh{1}$. The
\emph{half-open} convention, ranges $[l, r)$ that include $l$ and exclude $r$, is what makes the formula clean: there is no
$\pm 1$ anywhere, the empty range $l = r$ has sum $P_l - P_l = 0$, and $P$ has one more entry than $a$.

\begin{figure}[htbp]
\centering
\begin{tikzpicture}
  \node[blbl, anchor=east] at (-0.25, 0) {$a$};
  \arrs{A}{0}{0}{3/n, 1/n, 4/hot, 1/hot, 5/hot, 9/n, 2/n}
  \idx{A}{7}
  \node[blbl, anchor=east] at (-0.6, -1.6) {$P$};
  \arrs{P}{-0.35}{-1.6}{0/n, 3/n, 4/cool, 8/n, 9/n, 14/cool, 23/n, 25/n}
  \idxfrom{P}{8}{0}
  \spanabove{A-3}{A-5}{$a_2 + a_3 + a_4 = 10$}
  \node[notel, anchor=west] at (5.6, -1.6) {$= P_5 - P_2 = 14 - 4$};
\end{tikzpicture}
\caption{Prefix sums sit \emph{between} the elements: $P_i$ is the sum of everything to the left of boundary $i$. A range sum
is the difference of the prefix sums at its two boundaries.}
\label{fig:prefix}
\end{figure}

\usetemplate{Prefix sums (1D)}{Ctrl+Alt+T, 1}{cpprefix}

The same idea works for any operation that can be \emph{undone}. Prefix XORs give range XORs, since $x \oplus x = 0$. Prefix
counts give ``how many times does the letter \texttt{c} occur in $[l, r)$?'' for every letter (26 prefix arrays, or one array
of 26 counters per position). Prefix sums modulo a prime give range sums modulo that prime, and prefix \emph{products} give range
products if every element is invertible. Operations that cannot be undone, such as minimum, maximum and gcd, have no prefix
trick: a range minimum is not a difference of two prefix minima. Those need the sparse table or segment tree of
Part~VI.

\begin{trick}[Prefix sums without special cases]
Use the half-open convention everywhere and size $P$ as $n + 1$. When the input is 1-indexed and inclusive, convert once at
reading time (\cc{l--}, keep \cc{r}), not inside every formula. For circular arrays, append the array to itself and use
ranges of length at most $n$. And when values or counts may exceed $2^{31}$, the prefix sums certainly will: make them
\cc{long long} even if the elements fit in \cc{int}.
\end{trick}

\subsection{Two dimensions}

Let $P_{i,j}$ be the sum of the rectangle of rows $[0, i)$ and columns $[0, j)$. The sum of a rectangle with rows $[r_1, r_2)$
and columns $[c_1, c_2)$ follows by inclusion--exclusion: take everything above and to the left of its lower-right corner,
remove the strip above it and the strip to its left, and add back the corner region, which was removed twice
(Figure~\ref{fig:prefix2d}):
\[
  \mathrm{sum} = P_{r_2, c_2} - P_{r_1, c_2} - P_{r_2, c_1} + P_{r_1, c_1} .
\]
The table itself is built with the same identity in reverse, $P_{i+1, j+1} = P_{i, j+1} + P_{i+1, j} - P_{i,j} + a_{i,j}$, in
$\Oh{nm}$.

\begin{figure}[htbp]
\centering
\begin{tikzpicture}[x=0.26cm, y=0.26cm]
  \tikzset{qr/.style={draw=fgverm, line width=1.1pt}}
  \foreach \k/\op in {0/{$=$}, 1/{$-$}, 2/{$-$}, 3/{$+$}, 4/{}} {
    \begin{scope}[xshift=\k*3.1cm]
      \ifnum\k=0 \fill[fgverm!25] (3, -2) rectangle (7, -5);
      \else
        \ifnum\k=1 \fill[fgblue!25] (0, 0) rectangle (7, -5); \fi
        \ifnum\k=2 \fill[fgorange!30] (0, 0) rectangle (7, -2); \fi
        \ifnum\k=3 \fill[fgorange!30] (0, 0) rectangle (3, -5); \fi
        \ifnum\k=4 \fill[fggreen!35] (0, 0) rectangle (3, -2); \fi
      \fi
      \draw[fgline!70, line width=0.3pt, step=1] (0, 0) grid (8, -6);
      \draw[qr] (3, -2) rectangle (7, -5);
      \ifnum\k<4 \node[font=\large] at (9.5, -3) {\op}; \fi
    \end{scope}
  }
  \node[lbl] at (4, -7.2) {query};
  \node[lbl] at (4 + 11.92, -7.2) {$P_{r_2, c_2}$};
  \node[lbl] at (4 + 2*11.92, -7.2) {$P_{r_1, c_2}$};
  \node[lbl] at (4 + 3*11.92, -7.2) {$P_{r_2, c_1}$};
  \node[lbl] at (4 + 4*11.92, -7.2) {$P_{r_1, c_1}$};
\end{tikzpicture}
\caption{Two-dimensional prefix sums as a picture equation. The query rectangle (red frame) equals the region above and left
of its lower-right corner, minus the strip above it, minus the strip to its left, plus the corner region that both strips
removed.}
\label{fig:prefix2d}
\end{figure}

\begin{problem}{Forest Queries (CSES)}
An $n \times n$ grid ($n \le 1000$) contains trees. Answer $2\cdot10^5$ queries: how many trees are in a given rectangle?
\solution Build $P$ once, answer each query with the four-term formula. Converting the 1-indexed, inclusive rectangle to the
half-open form is the only place where care is needed.
\end{problem}

\begin{code}
void solve() {
    int n, q;
    cin >> n >> q;
    // P[i][j] = number of trees in rows < i and columns < j
    vvi P(n + 1, vi(n + 1, 0));
    rep(i, n) {
        string row;
        cin >> row;
        rep(j, n) {
            int tree = row[j] == '*';
            P[i + 1][j + 1] = P[i][j + 1] + P[i + 1][j] - P[i][j] + tree;
        }
    }
    while (q--) {
        int y1, x1, y2, x2;  // rows y1..y2, columns x1..x2, 1-indexed
        cin >> y1 >> x1 >> y2 >> x2;
        y1--, x1--;  // to the half-open rectangle [y1, y2) x [x1, x2)
        print(P[y2][x2] - P[y1][x2] - P[y2][x1] + P[y1][x1]);
    }
}
\end{code}

\usetemplate{2D prefix sums}{Ctrl+Alt+T, Shift+A}{cpprefix2d}

In $d$ dimensions the formula has $2^d$ terms with alternating signs, one per corner of the box. A different way to see
it, which generalises better, is to take prefix sums along one dimension at a time: first along the rows, then along the
columns of the result. After both passes, each entry is the sum of its whole upper-left rectangle. The same ``one dimension
at a time'' idea computes sums over all submasks of every bitmask (Chapter~\ref{ch:dppatterns}), which is prefix sums in
$n$ dimensions of size 2.

% ======================================================================================================================
\section[Difference arrays]{Difference arrays \hfill \silver}\label{sec:diff}

Prefix sums answer many queries about a fixed array. The inverse problem, applying many \emph{updates} to an array and reading
it only at the end, is solved by the inverse operation. The \defn{difference array} of $a$ is $D_i = a_i - a_{i-1}$ (with
$a_{-1} = 0$), so $a$ is the prefix sum of $D$. Adding $v$ to every element of $[l, r)$ changes only two differences:
$D_l$ grows by $v$ and $D_r$ shrinks by $v$. So each range update costs $\Oh{1}$, and one prefix-sum pass at the end rebuilds
the array (Figure~\ref{fig:diff}).

\begin{figure}[htbp]
\centering
\begin{tikzpicture}
  \node[blbl, anchor=east] at (-0.25, 0.9) {updates};
  \draw[ivl=fgblue!55] (0.7, 1.15) -- (2.8, 1.15);
  \node[font=\sffamily\scriptsize, text=fgblue, anchor=west] at (2.9, 1.15) {$+5$ on $[1, 4)$};
  \draw[ivl=fgorange!80] (2.1, 0.65) -- (4.2, 0.65);
  \node[font=\sffamily\scriptsize, text=fgorange!80!black, anchor=west] at (4.3, 0.65) {$+2$ on $[3, 6)$};
  \node[blbl, anchor=east] at (-0.25, -0.4) {$D$};
  \arrs{D}{0}{-0.4}{0/n, +5/cool, 0/n, +2/warn, $-5$/cool, 0/n, $-2$/warn}
  \idx{D}{7}
  \node[blbl, anchor=east] at (-0.25, -1.9) {prefix sums};
  \arrs{S}{0}{-1.9}{0/n, 5/good, 5/good, 7/good, 2/good, 2/good, 0/n}
  \node[notel, anchor=west, text width=5.2cm] at (5.4, -1.15) {each update writes two cells of $D$; one pass of
    prefix sums applies all of them at once};
\end{tikzpicture}
\caption{A difference array. The update $+5$ on $[1, 4)$ becomes $+5$ at index 1 and $-5$ at index 4; after all updates, the
prefix sums of $D$ are the final values.}
\label{fig:diff}
\end{figure}

\usetemplate{Difference array (offline range add, 1D)}{Ctrl+Alt+T, 3}{cpdiff}

\subsection{Rectangles}

In two dimensions, adding $v$ to the rectangle $[x_1, x_2) \times [y_1, y_2)$ writes four corners of the difference array:
$+v$ at $(x_1, y_1)$ and $(x_2, y_2)$, and $-v$ at $(x_1, y_2)$ and $(x_2, y_1)$. After 2D prefix sums, the $+v$ at
$(x_1, y_1)$ spreads to the whole quadrant below and to the right of it; the two $-v$ remove the parts beyond $x_2$ and
beyond $y_2$; and the second $+v$ restores the part beyond both, which was removed twice (Figure~\ref{fig:diff2d}). This is
inclusion--exclusion again, run backwards.

\begin{figure}[htbp]
\centering
\begin{tikzpicture}[x=0.42cm, y=0.42cm]
  \begin{scope}
    \draw[fgline!70, line width=0.3pt, step=1] (0, 0) grid (8, -6);
    \foreach \x/\y/\s/\c in {2/1/+1/fgblue, 6/1/$-1$/fgverm, 2/4/$-1$/fgverm, 6/4/+1/fgblue}
      \node[font=\sffamily\scriptsize\bfseries, text=\c] at (\x + 0.5, -\y - 0.5) {\s};
    \node[lbl] at (4, -6.8) {difference array: four corners};
  \end{scope}
  \node[font=\large] at (10, -3) {$\Rightarrow$};
  \begin{scope}[xshift=5.3cm]
    \fill[fggreen!30] (2, -1) rectangle (6, -4);
    \draw[fgline!70, line width=0.3pt, step=1] (0, 0) grid (8, -6);
    \draw[line width=1.1pt, draw=fggreen!70!black] (2, -1) rectangle (6, -4);
    \foreach \x in {2,...,5} \foreach \y in {1,...,3}
      \node[font=\sffamily\scriptsize, text=fggreen!50!black] at (\x + 0.5, -\y - 0.5) {1};
    \node[lbl] at (4, -6.8) {after 2D prefix sums};
  \end{scope}
\end{tikzpicture}
\caption{Adding 1 to the rectangle of columns 2--5 and rows 1--3. The difference array gets $+1$ at the rectangle's top-left
cell and at the cell diagonally beyond its bottom-right corner, and $-1$ at the cells just beyond its other two corners; 2D
prefix sums turn these four values into exactly the rectangle.}
\label{fig:diff2d}
\end{figure}

\begin{problem}{Painting the Barn (USACO Silver)}
Farmer John paints $n \le 10^5$ axis-parallel rectangles on a barn wall with coordinates in $[0, 1000]$; each rectangle adds
one coat of paint to the unit squares it covers. How many unit squares end up with exactly $k$ coats?
\solution Four updates per rectangle into a $1001 \times 1001$ difference array, one 2D prefix-sum pass, and a count: about
$10^6$ operations, regardless of how large the rectangles are.
\end{problem}

\begin{code}
void solve() {
    int n, k;
    cin >> n >> k;
    const int N = 1001;
    vvi D(N + 1, vi(N + 1, 0));  // 2D difference array
    rep(i, n) {
        int x1, y1, x2, y2;  // the unit squares [x1, x2) x [y1, y2)
        cin >> x1 >> y1 >> x2 >> y2;
        D[x1][y1]++, D[x1][y2]--, D[x2][y1]--, D[x2][y2]++;
    }
    int cells = 0;
    rep(x, N) rep(y, N) {  // 2D prefix sums in place: D[x][y] = coats on (x, y)
        if (x > 0) D[x][y] += D[x - 1][y];
        if (y > 0) D[x][y] += D[x][y - 1];
        if (x > 0 && y > 0) D[x][y] -= D[x - 1][y - 1];
        if (D[x][y] == k) cells++;
    }
    print(cells);
}
\end{code}

\usetemplate{Difference array (offline rectangle add, 2D)}{Ctrl+Alt+T, 4}{cpdiff2d}

\subsection{Arithmetic progressions: differences of differences}

What if an update adds $1, 2, 3, \ldots$ to the elements of a range? The first differences of such an update are not two
spikes but a whole block of ones, so one level of differences is not enough. Take differences \emph{twice}: the second
differences of the progression are $+1$ at its start, $-(len + 1)$ just after its end and $+len$ one step later, three spikes
again (Figure~\ref{fig:diff-ap}). Two prefix-sum passes rebuild the array. In general, an update by a polynomial of degree
$d$ becomes $\Oh{d}$ spikes after $d + 1$ levels of differences, and $d + 1$ prefix-sum passes undo them.

\begin{figure}[htbp]
\centering
\begin{tikzpicture}
  \node[blbl, anchor=east] at (-0.25, 0) {second differences};
  \arrs{E}{0}{0}{0/n, 0/n, +1/hot, 0/n, 0/n, 0/n, $-5$/hot, +4/hot}
  \node[blbl, anchor=east] at (-0.25, -1.1) {first differences};
  \arrs{F}{0}{-1.1}{0/n, 0/n, 1/cool, 1/cool, 1/cool, 1/cool, $-4$/cool, 0/n}
  \node[blbl, anchor=east] at (-0.25, -2.2) {values};
  \arrs{V}{0}{-2.2}{0/n, 0/n, 1/good, 2/good, 3/good, 4/good, 0/n, 0/n}
  \idx{V}{8}
  \node[notel, anchor=west, text width=4.6cm] at (6.0, -1.1) {adding $1, 2, 3, 4$ to positions $2..5$ ($len = 4$):
    three spikes in the second differences; two passes of prefix sums rebuild the values};
\end{tikzpicture}
\caption{An arithmetic progression added to a range is three spikes in the second differences.}
\label{fig:diff-ap}
\end{figure}

\begin{code}
void solve() {
    int n, q;
    cin >> n >> q;
    vi d2(n + 2, 0);  // second differences
    while (q--) {
        int l, r;  // add 1, 2, ..., r - l + 1 to positions l..r (0-indexed)
        cin >> l >> r;
        int len = r - l + 1;
        d2[l] += 1;  // the first differences become 1 on [l, r]
        d2[r + 1] -= len + 1;  // ... then drop by len right after r
        d2[r + 2] += len;  // ... and are back to 0 after that
    }
    vi a(n);
    ll d1 = 0, v = 0;
    rep(i, n) {
        d1 += d2[i];  // first difference a[i] - a[i - 1]
        v += d1;
        a[i] = v;
    }
    print(a);
}
\end{code}

\begin{recognize}
Many updates, few (or one) reads, and the updates add something simple to a contiguous range: an interval of days, a
rectangle of cells, a path in a tree (Chapter~\ref{ch:trees}). When reads and updates are interleaved, the same two-point
update goes into a Fenwick tree instead, which supports prefix sums online (Chapter~\ref{ch:segtree}).
\end{recognize}

% ======================================================================================================================
\section[Counting subarrays with a hash map]{Counting subarrays with a hash map \hfill \silver}\label{sec:prefix-hash}

``How many subarrays have sum exactly $x$?'' With prefix sums, the subarray $[l, r)$ has sum $x$ exactly when $P_r - P_l = x$,
that is, when $P_l = P_r - x$. So scan $r$ from left to right, keep a map from each prefix sum seen so far to the number of
times it has occurred, and add the count of $P_r - x$ at each step: every earlier prefix with that value is the left end of a
subarray ending at $r$ with sum $x$. This counts all subarrays in $\Oh{n \log n}$ with a \cc{map} (or expected $\Oh{n}$ with
a hash map), even when the values are negative, which rules out two pointers. The complete program is the worked problem
Subarray Sums II (Chapter~\ref{ch:worked}).

Picturing the prefix sums as a walk makes these counting problems visual: the subarrays with sum 0 are exactly the pairs of
moments at which the walk is at the same height (Figure~\ref{fig:zerosum}), and the subarrays with sum $x$ are the pairs of
moments whose heights differ by $x$.

\begin{figure}[htbp]
\centering
\begin{tikzpicture}[x=1.0cm, y=0.55cm]
  \draw[->, line width=0.5pt, draw=fgink!60] (-0.3, 0) -- (7.6, 0) node[right, lbl] {$i$};
  \foreach \i in {0,...,7} \node[lbl] at (\i, -0.55) {\i};
  \foreach \h in {0, 2} \draw[line width=0.5pt, draw=fgline, dashed] (-0.2, \h) -- (7.3, \h);
  \draw[line width=1.3pt, draw=fgblue] (0, 0) -- (1, 3) -- (2, 1) -- (3, 2) -- (4, 0) -- (5, 4) -- (6, 0) -- (7, 2);
  \foreach \i/\p in {0/0, 1/3, 2/1, 3/2, 4/0, 5/4, 6/0, 7/2} {
    \fill[fgblue] (\i, \p) circle (2.2pt);
    \node[font=\sffamily\scriptsize, text=fgblue, anchor=south] at (\i, \p + 0.15) {\p};
  }
  \foreach \a/\b/\yy/\cl in {0/4/-1.4/fggreen, 4/6/-1.4/fggreen, 0/6/-2.1/fggreen, 3/7/-2.8/fgorange}
    \draw[ivl=\cl!75] (\a + 0.06, \yy) -- (\b - 0.06, \yy);
  \node[notel, anchor=north west, text width=4.3cm] at (7.9, 4.4) {$a = 3, -2, 1, -2, 4, -4, 2$\\ prefix sums $P_0, \ldots,
    P_7$ (blue); equal heights are zero-sum subarrays: $[0, 4)$, $[4, 6)$, $[0, 6)$ at height 0 and $[3, 7)$ at height 2};
\end{tikzpicture}
\caption{Prefix sums as a walk. Every pair of equal heights is a subarray with sum 0; the bars below the axis show the four
such subarrays (green at height 0, orange at height 2). Counting them is counting pairs of equal values in $P$.}
\label{fig:zerosum}
\end{figure}

The same scan answers a family of questions; only the map and its lookup change.
\begin{itemize}
  \item \textbf{Sum divisible by $m$:} two prefixes with equal $P \bmod m$ (Section~\ref{sec:pigeonhole}); normalise negative
        remainders with \cc{((P % m) + m) % m}.
  \item \textbf{XOR equal to $x$:} prefix XORs, and $X_r \oplus X_l = x \iff X_l = X_r \oplus x$.
  \item \textbf{Equal numbers of two kinds of elements:} count one kind as $+1$ and the other as $-1$, and look for zero sums.
  \item \textbf{The longest such subarray:} store the \emph{first} index of each prefix value instead of a count; the
        longest subarray ending at $r$ starts at the first occurrence of the value it needs.
\end{itemize}

\begin{problem}{Longest balanced substring}
Find the length of the longest substring of a binary string that contains equally many zeros and ones.
\solution With $+1$ for each one and $-1$ for each zero, balanced substrings are exactly the zero-sum subarrays. The longest
one ending at position $i$ starts right after the \emph{first} prefix with the same sum, so a map from prefix sum to the first
index where it occurred gives the answer in one pass.
\end{problem}

\begin{code}
void solve() {
    int n;
    cin >> n;
    string s;
    cin >> s;
    // count 1 as +1 and 0 as -1: balanced means a zero sum
    map<int, int> first;  // the first index at which each prefix sum occurs
    first[0] = 0;
    int pre = 0, best = 0;
    rep(i, n) {
        pre += s[i] == '1' ? 1 : -1;
        if (first.count(pre)) chmax(best, i + 1 - first[pre]);
        else first[pre] = i + 1;
    }
    print(best);
}
\end{code}

\begin{pitfall}[Hash maps and adversarial tests]
An \cc{unordered_map} of prefix sums is expected $\Oh{1}$ per operation, but on Codeforces, tests are often built to collide
with the default hash, which makes every operation linear. Use a \cc{map}, a custom hash with a random seed
(Chapter~\ref{ch:cpp}), or, when the prefix values lie in a small range (as in the balanced-string problem, where they lie in
$[-n, n]$), a plain array with an offset.
\end{pitfall}

% ======================================================================================================================
\section[Two pointers]{Two pointers \hfill \silver}\label{sec:two-pointers}

Suppose we want, for every right end $r$, the smallest left end $l$ such that the window $[l, r]$ satisfies some condition,
and the condition is \emph{monotone}: shrinking a valid window keeps it valid (for example ``the sum is at most $x$'' with
non-negative values, or ``all values are distinct''). Then the best left end $l(r)$ never moves backwards when $r$ grows,
because a window that was invalid stays invalid when it is extended to the right. So the two ends can be moved with two
pointers, both only forwards: for each new $r$, advance $l$ while the window is invalid.

\emph{Why it is linear.} Each step of the algorithm moves one pointer one position forwards, and each pointer moves at most
$n$ times, so there are at most $2n$ steps in total, however the work is distributed among the iterations. Plotting the pair
$(r, l)$ as a point, the algorithm walks a monotone staircase path from the bottom-left to the top-right corner of an
$n \times n$ grid (Figure~\ref{fig:staircase}). This is the simplest instance of amortised analysis
(Section~\ref{sec:amortized}).

\begin{figure}[htbp]
\centering
\begin{tikzpicture}[x=0.62cm, y=0.62cm]
  \fill[fgline!22] (0, 1) \foreach \k in {1,...,6} { -- (\k, \k) -- (\k, \k + 1) } -- (0, 7) -- cycle;
  \draw[fgline!60, line width=0.3pt, step=1] (0, 0) grid (7, 7);
  \draw[->, line width=0.6pt, draw=fgink!70] (0, 0) -- (7.6, 0) node[right, lbl] {$r$};
  \draw[->, line width=0.6pt, draw=fgink!70] (0, 0) -- (0, 7.6) node[above, lbl] {$l$};
  \foreach \i in {0,...,6} { \node[lbl] at (\i + 0.5, -0.4) {\i}; \node[lbl] at (-0.4, \i + 0.5) {\i}; }
  \draw[line width=1.6pt, draw=fgblue] (0.5, 0.5) -- (2.5, 0.5) -- (3.5, 0.5) -- (3.5, 1.5) -- (4.5, 1.5) -- (4.5, 4.5)
    -- (5.5, 4.5) -- (5.5, 5.5) -- (6.5, 5.5);
  \foreach \r/\l/\ok in {0/0/0, 1/0/0, 2/0/1, 3/1/1, 4/4/1, 5/5/0, 6/5/0} {
    \ifnum\ok=1 \fill[fggreen] (\r + 0.5, \l + 0.5) circle (4pt); \else \fill[fgblue] (\r + 0.5, \l + 0.5) circle (2.6pt); \fi
  }
  \node[notel, anchor=west, text width=5.6cm] at (8.2, 4.2) {Subarray Sums I with $a = 2, 4, 1, 2, 7, 1, 3$ and $x = 7$.
    For each $r$ the dot marks the left end $l(r)$ after shrinking; green dots are windows with sum exactly 7:
    $[0, 2]$, $[1, 3]$ and $[4, 4]$. The path only goes right or up, so it has at most $2n$ steps.};
\end{tikzpicture}
\caption{Two pointers as a staircase. Moving $r$ is a step to the right, moving $l$ a step up; the shaded region ($l > r$)
is never entered here. The total number of steps is at most $2n$, although a single value of $r$ can move $l$ many times.}
\label{fig:staircase}
\end{figure}

\begin{problem}{Subarray Sums I (CSES)}
Count the subarrays of an array of $n \le 2\cdot10^5$ \emph{positive} integers whose sum is exactly $x$.
\solution With positive values, the sum of a window only grows when it is extended and only shrinks when it is cut, so for
each right end there is at most one left end with sum $x$, and it can be found with two pointers: add $a_r$, then remove
elements from the left while the sum exceeds $x$.
\end{problem}

\begin{code}
void solve() {
    int n;
    ll x;
    cin >> n >> x;
    vi a(n);
    cin >> a;
    ll sum = 0, ways = 0;  // all values are positive
    int l = 0;
    rep(r, n) {  // the window is a[l..r]
        sum += a[r];
        while (sum > x) sum -= a[l++];  // too large: shrink from the left
        if (sum == x) ways++;
    }
    print(ways);
}
\end{code}

\usetemplate{Two pointers (sliding window)}{Ctrl+Alt+T, 5}{cptwoptr}

\begin{problem}{Playlist (CSES)}
Find the longest contiguous part of a playlist in which no song repeats.
\solution ``All songs distinct'' is monotone under shrinking. For each right end, the left end must move past the previous
occurrence of the new song if that occurrence is inside the window. Storing the last position of every song lets the left
pointer jump there directly.
\end{problem}

\begin{code}
void solve() {
    int n;
    cin >> n;
    vi k(n);
    cin >> k;
    map<ll, int> last;  // the last position of each song so far
    int l = 0, best = 0;
    rep(r, n) {  // the longest window of distinct songs that ends at r
        auto it = last.find(k[r]);
        if (it != last.end() && it->second >= l) l = it->second + 1;
        last[k[r]] = r;
        chmax(best, r - l + 1);
    }
    print(best);
}
\end{code}

\subsection{Counting windows: at most, and exactly}

When the condition is monotone, all windows ending at $r$ with left end at least $l(r)$ are valid, so the number of valid
windows is $\sum_r (r - l(r) + 1)$, again in linear time.

\begin{problem}{Subarray Distinct Values (CSES)}
Count the subarrays that contain at most $k$ distinct values.
\solution Keep the count of each value in the window in a map; the number of distinct values is the size of the map. For
each $r$, add $a_r$ and shrink from the left while there are more than $k$ distinct values; then add $r - l + 1$.
\end{problem}

\begin{code}
void solve() {
    int n, k;
    cin >> n >> k;
    vi a(n);
    cin >> a;
    map<ll, int> cnt;  // occurrences of each value in the window a[l..r]
    ll ways = 0;
    int l = 0;
    rep(r, n) {
        cnt[a[r]]++;
        while (sz(cnt) > k) {  // too many distinct values: shrink
            if (--cnt[a[l]] == 0) cnt.erase(a[l]);
            l++;
        }
        ways += r - l + 1;  // every window a[l'..r] with l' >= l is fine
    }
    print(ways);
}
\end{code}

\begin{trick}[Exactly $k$ = at most $k$ $-$ at most $k - 1$]
``Exactly $k$ distinct values'' is not monotone: shrinking a window can drop it to $k - 1$. But
\[
  \#\{\text{exactly } k\} = \#\{\text{at most } k\} - \#\{\text{at most } k - 1\},
\]
and both terms on the right are monotone conditions that two pointers count directly. The same subtraction handles ``sum
exactly $x$'' with non-negative values (at most $x$ minus at most $x - 1$), ``exactly $k$ odd numbers'' and many more.
\end{trick}

\subsection{Pointers from both ends}

The other classic form of two pointers starts one pointer at each end of a sorted array and moves them towards each other.

\begin{problem}{Sum of Two Values (CSES)}
Find two positions (distinct) whose values add up to $x$, or report that there are none.
\solution Sort the values (remembering their original positions) and start with $i$ at the smallest and $j$ at the largest.
If $a_i + a_j < x$, then $a_i$ is too small even for the largest remaining partner, so it can be discarded: $i$ moves right.
If $a_i + a_j > x$, then $a_j$ is too large even for the smallest remaining partner: $j$ moves left. Each step discards one
value for good, which is why the scan is linear and misses nothing (Figure~\ref{fig:saddleback}).
\end{problem}

\begin{figure}[htbp]
\centering
\begin{tikzpicture}[x=0.72cm, y=0.62cm]
  \foreach \i/\vi in {0/1, 1/3, 2/4, 3/6, 4/8} \foreach \j/\vj in {1/3, 2/4, 3/6, 4/8, 5/11} {
    \pgfmathtruncatemacro{\s}{\vi + \vj}
    \ifnum\j>\i
      \node[draw=fgline, line width=0.4pt, fill=white, minimum width=7mm, minimum height=5.8mm, inner sep=0pt,
        font=\ttfamily\small] (m\i\j) at (\j, -\i) {\s};
    \fi
  }
  \foreach \i/\j in {0/5, 0/4, 1/4, 1/3} \node[draw=fgblue, fill=fgsky!25, line width=0.8pt, minimum width=7mm,
    minimum height=5.8mm, inner sep=0pt, font=\ttfamily\small] at (\j, -\i) {};
  \foreach \i/\j/\s in {0/5/12, 0/4/9, 1/4/11, 1/3/9} \node[font=\ttfamily\small] at (\j, -\i) {\s};
  \node[draw=fggreen, fill=fggreen!25, line width=1.2pt, minimum width=7mm, minimum height=5.8mm, inner sep=0pt,
    font=\ttfamily\small\bfseries] at (3, -2) {10};
  \draw[line width=5pt, draw=fgverm, opacity=0.28, line join=round] (5, 0) -- (4, 0) -- (4, -1) -- (3, -1)
    -- (3, -2);
  \foreach \j/\vj in {1/3, 2/4, 3/6, 4/8, 5/11} \node[blbl] at (\j, 0.85) {\vj};
  \foreach \i/\vi in {0/1, 1/3, 2/4, 3/6, 4/8} \node[blbl] at (-0.1, -\i) {\vi};
  \node[lbl] at (3, 1.5) {$a_j$};
  \node[lbl, anchor=east] at (-0.55, -2) {$a_i$};
  \node[notel, anchor=west, text width=5.8cm] at (6.1, -1.8) {Sums $a_i + a_j$ for $a = 1, 3, 4, 6, 8, 11$ and target 10.
    They increase downwards and to the right. From the top-right corner, a sum that is too large rules out its whole
    column (move left), and one that is too small rules out its whole row (move down).};
\end{tikzpicture}
\caption{Two pointers from both ends as a walk through the table of pair sums: $12 > 10$ (left), $9 < 10$ (down), $11 > 10$
(left), $9 < 10$ (down), $10$ found. Each step discards a whole row or column, so at most $n$ steps are needed.}
\label{fig:saddleback}
\end{figure}

\begin{code}
void solve() {
    int n;
    ll x;
    cin >> n >> x;
    vector<pll> a(n);  // (value, original position)
    rep(i, n) cin >> a[i].first, a[i].second = i + 1;
    sort(all(a));
    int i = 0, j = n - 1;
    while (i < j) {
        ll s = a[i].first + a[j].first;
        if (s == x) {
            print(a[i].second, a[j].second);
            return;
        }
        if (s < x) i++;  // a[i] is too small for every remaining partner
        else j--;  // a[j] is too large for every remaining partner
    }
    print("IMPOSSIBLE");
}
\end{code}

Fixing one element and running this scan on the rest solves 3-SUM (Sum of Three Values) in $\Oh{n^2}$; for 4-SUM, fix two
elements, or use meet in the middle over pairs (Section~\ref{sec:mitm}).

\begin{recognize}
``The longest or shortest subarray such that\ldots'', ``count the subarrays with at most\ldots'', pairs in a sorted array with
a given sum or difference. The key test is monotonicity: does shrinking a valid window keep it valid? With non-negative
values sums are monotone; with negative values they are not, and prefix sums with a map or a data structure take over.
\end{recognize}

% ======================================================================================================================
\section[Sliding windows]{Sliding windows \hfill \gold}\label{sec:sliding}

When every window has the same length $k$, the question is usually an aggregate of the window: its minimum, maximum, sum,
median or number of distinct values. The sum and the number of distinct values are updated in $\Oh{1}$ as one element enters
and one leaves. The minimum and maximum need the monotonic deque of Section~\ref{sec:queues}, which is amortised $\Oh{1}$ per
step; for any other associative operation, the two-stack queue of Section~\ref{sec:augment} works. The median is harder,
because it needs order statistics of a multiset that changes at both ends.

\begin{problem}{Sliding Window Median (CSES)}
For each window of $k$ consecutive elements, print its median (the lower one when $k$ is even).
\solution Split the window into two multisets: \emph{low} holds the $\lceil k/2 \rceil$ smallest values and \emph{high} the
rest, so the median is the largest element of \emph{low} (Figure~\ref{fig:slidingmedian}). A new element goes into the set
it belongs to by comparison with that median; the leaving element is erased from whichever set contains it; then at most one
element moves between the sets to restore the sizes. Every step is $\Oh{\log k}$.
\end{problem}

\begin{figure}[htbp]
\centering
\begin{tikzpicture}
  \foreach \row/\wa/\wb/\low/\high/\med/\note in {0/1/5/{1,3,5}/{8,9}/5/{window $5, 1, 8, 3, 9$},
      1/2/6/{1,2,3}/{8,9}/3/{5 leaves, 2 enters}} {
    \pgfmathsetmacro{\yy}{-1.55*\row}
    \foreach \v [count=\i] in {5, 1, 8, 3, 9, 2, 7} {
      \pgfmathtruncatemacro{\inw}{(\i >= \wa && \i <= \wb) ? 1 : 0}
      \ifnum\inw=1 \node[cellc, cool] at (\i*0.7, \yy) {\v}; \else \node[cellc, dim] at (\i*0.7, \yy) {\v}; \fi
    }
    \node[lbl, anchor=east] at (0.25, \yy) {\note};
    \node[draw=fgblue, fill=fgsky!20, rounded corners=2pt, line width=0.7pt, minimum height=6mm, inner sep=3pt,
      font=\sffamily\small] (lo\row) at (7.3, \yy) {low: \low};
    \node[draw=fgorange!85!black, fill=fgorange!15, rounded corners=2pt, line width=0.7pt, minimum height=6mm,
      inner sep=3pt, font=\sffamily\small, anchor=west] at ($(lo\row.east)+(0.25, 0)$) {high: \high};
    \node[font=\sffamily\small\bfseries, text=fgverm, anchor=west] at (10.6, \yy) {median \med};
  }
\end{tikzpicture}
\caption{Sliding window median with $k = 5$. \emph{low} keeps the three smallest values of the window, \emph{high} the other
two; the median is the largest value in \emph{low}. When the window slides, 5 is erased and 2 inserted, and the median
becomes 3.}
\label{fig:slidingmedian}
\end{figure}

\begin{code}
void solve() {
    int n, k;
    cin >> n >> k;
    vi a(n);
    cin >> a;
    // low: the (k + 1) / 2 smallest values of the window, high: the rest
    multiset<ll> low, high;
    auto balance = [&]() {
        while (sz(low) > (sz(low) + sz(high) + 1) / 2) {
            high.insert(*prev(low.end()));
            low.erase(prev(low.end()));
        }
        while (sz(low) < (sz(low) + sz(high) + 1) / 2) {
            low.insert(*high.begin());
            high.erase(high.begin());
        }
    };
    vi med;
    rep(i, n) {
        if (low.empty() || a[i] <= *prev(low.end())) low.insert(a[i]);
        else high.insert(a[i]);
        if (i >= k) {  // a[i - k] leaves the window
            auto it = low.find(a[i - k]);
            if (it != low.end()) low.erase(it);
            else high.erase(high.find(a[i - k]));
        }
        balance();
        if (i >= k - 1) med.push_back(*prev(low.end()));  // largest of low
    }
    print(med);
}
\end{code}

Two details deserve attention. Erasing a value with \cc{low.erase(x)} would remove \emph{all} copies of it; \cc{erase} of an
iterator removes exactly one. And the leaving element must be erased from the set that actually contains it, which the code
decides by searching \emph{low} first. Keeping the sums of both sets as well gives Sliding Window Cost, the minimum total
distance from the window's values to a single point: by Section~\ref{sec:greedy-classics} that point is the median, and the
cost is $(\text{med} \cdot |low| - \Sigma_{low}) + (\Sigma_{high} - \text{med} \cdot |high|)$.

% ======================================================================================================================
\section[Monotonic stacks in action]{Monotonic stacks in action \hfill \gold}\label{sec:monostack-apps}

The monotonic stack of Section~\ref{sec:stacks} computes, for every element, the nearest smaller (or larger) element on each
side in $\Oh{n}$. Those two numbers describe the largest window in which an element is the minimum, and that is exactly what
many window problems need: the contribution technique of Section~\ref{sec:contribution} counts the subarrays in which each
element is the minimum, and the following classic finds the largest rectangle under a histogram.

\begin{problem}{Advertisement (CSES)}
A fence consists of $n$ boards of width 1 and heights $h_i$. Find the largest rectangular advertisement that can be placed on
the fence, aligned with the boards.
\solution The largest rectangle has the height of its lowest board $i$ (otherwise it could be made taller), and extends left
and right until a lower board blocks it. With $L_i$ and $R_i$, the nearest lower boards on each side from two monotonic stack
passes, the rectangle for board $i$ has width $R_i - L_i - 1$, and the answer is $\max_i h_i (R_i - L_i - 1)$
(Figure~\ref{fig:histogram}).
\end{problem}

\begin{figure}[htbp]
\centering
\begin{tikzpicture}[x=0.8cm, y=0.45cm]
  \fill[fggreen!30] (2, 0) rectangle (4, 5);
  \foreach \i/\h in {0/2, 1/1, 2/5, 3/6, 4/2, 5/3} {
    \ifnum\i=1 \def\bc{fgverm}\else\ifnum\i=4 \def\bc{fgverm}\else\def\bc{fgblue}\fi\fi
    \draw[line width=0.8pt, draw=\bc] (\i + 0.04, 0) rectangle (\i + 0.96, \h);
    \node[font=\sffamily\scriptsize] at (\i + 0.5, \h + 0.5) {\h};
    \node[lbl] at (\i + 0.5, -0.6) {\i};
  }
  \draw[line width=1.4pt, draw=fggreen!70!black] (2, 0) rectangle (4, 5);
  \node[font=\sffamily\scriptsize\bfseries, text=fggreen!40!black, fill=fggreen!30, inner sep=1.5pt] at (3, 2.5)
    {$5 \times 2 = 10$};
  \node[font=\sffamily\scriptsize\bfseries, text=fgverm] at (1.5, 2.3) {$L_2$};
  \node[font=\sffamily\scriptsize\bfseries, text=fgverm] at (4.5, 3.3) {$R_2$};
  \draw[<->, line width=0.9pt, draw=fggreen!60!black] (2, -1.5) -- node[below, font=\sffamily\scriptsize,
    text=fggreen!40!black] {width $R_2 - L_2 - 1 = 2$} (4, -1.5);
  \draw[line width=0.6pt, draw=fgink!70] (-0.2, 0) -- (6.2, 0);
\end{tikzpicture}
\caption{The largest rectangle in the histogram $2, 1, 5, 6, 2, 3$ has height 5, set by board 2. It extends from just after
the nearest lower board on the left ($L_2 = 1$) to just before the nearest lower board on the right ($R_2 = 4$); both
blockers are drawn in red.}
\label{fig:histogram}
\end{figure}

\begin{code}
void solve() {
    int n;
    cin >> n;
    vi h(n);
    cin >> h;
    // L[i], R[i]: nearest bars to the left and right lower than h[i]
    vi L(n), R(n), st;
    rep(i, n) {
        while (!st.empty() && h[st.back()] >= h[i]) st.pop_back();
        L[i] = st.empty() ? -1 : st.back();
        st.push_back(i);
    }
    st.clear();
    per(i, n) {
        while (!st.empty() && h[st.back()] >= h[i]) st.pop_back();
        R[i] = st.empty() ? n : st.back();
        st.push_back(i);
    }
    ll best = 0;
    rep(i, n) chmax(best, h[i] * (R[i] - L[i] - 1));  // as tall as bar i
    print(best);
}
\end{code}

\usetemplate{Monotonic stack (next greater / previous smaller)}{Ctrl+Alt+T, 7}{cpmonostack}

\paragraph{From one dimension to two.} The largest rectangle of free cells in an $n \times m$ grid (Maximum Building I)
reduces to $n$ histograms: for each row, let $h_j$ be the number of consecutive free cells ending at this row in column $j$
(zero if the cell is blocked, otherwise one more than in the previous row). The best rectangle whose bottom edge lies on this
row is the largest rectangle of that histogram, so the whole grid takes $\Oh{nm}$. Counting all free rectangles instead of
finding the largest one (Maximum Building II) uses the same heights with the contribution technique.

\begin{summary}[Chapter summary]
\begin{itemize}
  \item Prefix sums answer range sums in $\Oh{1}$ after $\Oh{n}$ work; use half-open ranges and $n + 1$ entries. In two
        dimensions, inclusion--exclusion over four corners.
  \item Difference arrays apply range updates in $\Oh{1}$ each and rebuild with one prefix-sum pass; rectangles need four
        corners, arithmetic progressions a second level of differences.
  \item Counting subarrays with sum $x$ is counting pairs of prefix sums with difference $x$: scan with a map. Variations
        handle remainders, XOR, balance, and the longest such subarray.
  \item Two pointers work when the window condition is monotone; both pointers move forwards, so the scan is linear. Count
        ``exactly'' as ``at most $k$'' minus ``at most $k-1$''.
  \item Two multisets maintain a sliding median; monotonic stacks give the largest window in which each element is the
        minimum, which solves the largest rectangle problems.
\end{itemize}
\end{summary}

\begin{practice}
\ppgroup{Prefix sums and differences}
\begin{pplist}
  \item Static Range Sum Queries
  \item Forest Queries
  \item Subarray Sums II
  \item Subarray Divisibility
\end{pplist}
\ppgroup{Two pointers}
\begin{pplist}
  \item Subarray Sums I
  \item Sum of Two Values
  \item Sum of Three Values
  \item Sum of Four Values
  \item Playlist
  \item Subarray Distinct Values
\end{pplist}
\ppgroup{Sliding windows and stacks}
\begin{pplist}
  \item Sliding Window Median
  \item Sliding Window Cost
  \item Nearest Smaller Values
  \item Advertisement
  \item Maximum Building I
\end{pplist}
\end{practice}
